\section{Sections over principal opens}
\label{section-principal-opens}

\noindent
Here is a typical result of this kind. We will use a more naive but
more direct method of proof in later lemmas.

\begin{lemma}
\label{lemma-invert-f-sections}
\begin{slogan}
Sections of quasi-coherent sheaves have only meromorphic singularities
at infinity.
\end{slogan}
Let $X$ be a scheme. Let $f \in \Gamma(X, \mathcal{O}_X)$.
Denote $X_f \subset X$ the open where $f$ is invertible, see
Schemes, Lemma \ref{schemes-lemma-f-open}.
If $X$ is quasi-compact and quasi-separated, the canonical map
$$
\Gamma(X, \mathcal{O}_X)_f \longrightarrow \Gamma(X_f, \mathcal{O}_X)
$$
is an isomorphism. Moreover, if $\mathcal{F}$ is a quasi-coherent
sheaf of $\mathcal{O}_X$-modules the map
$$
\Gamma(X, \mathcal{F})_f \longrightarrow \Gamma(X_f, \mathcal{F})
$$
is an isomorphism.
\end{lemma}

\begin{proof}
Write $R = \Gamma(X, \mathcal{O}_X)$. Consider the canonical morphism
$$
\varphi : X \longrightarrow \Spec(R)
$$
of schemes, see
Schemes, Lemma
\ref{schemes-lemma-morphism-into-affine}.
Then the inverse image of the standard open $D(f)$ on the
right hand side is $X_f$ on the left hand side.
Moreover, since $X$ is assumed quasi-compact and quasi-separated
the morphism $\varphi$ is quasi-compact and quasi-separated,
see Schemes, Lemma \ref{schemes-lemma-quasi-compact-affine} and
\ref{schemes-lemma-compose-after-separated}. Hence by
Schemes, Lemma \ref{schemes-lemma-push-forward-quasi-coherent}
we see that $\varphi_*\mathcal{F}$ is quasi-coherent.
Hence we see that $\varphi_*\mathcal{F} = \widetilde M$
with $M = \Gamma(X, \mathcal{F})$ as an $R$-module.
Thus we see that
$$
\Gamma(X_f, \mathcal{F}) =
\Gamma(D(f), \varphi_*\mathcal{F}) =
\Gamma(D(f), \widetilde M) = M_f
$$
which is exactly the content of the lemma. The first displayed isomorphism
of the lemma follows by taking $\mathcal{F} = \mathcal{O}_X$.
\end{proof}

\noindent
Recall that given a scheme $X$, an invertible sheaf $\mathcal{L}$
on $X$, and a sheaf of $\mathcal{O}_X$-modules $\mathcal{F}$
we get a graded ring
$\Gamma_*(X, \mathcal{L}) =
\bigoplus\nolimits_{n \geq 0} \Gamma(X, \mathcal{L}^{\otimes n})$
and a graded $\Gamma_*(X, \mathcal{L})$-module
$\Gamma_*(X, \mathcal{L}, \mathcal{F}) =
\bigoplus\nolimits_{n \in \mathbf{Z}}
\Gamma(X, \mathcal{F} \otimes_{\mathcal{O}_X} \mathcal{L}^{\otimes n})$
see Modules, Definition \ref{modules-definition-gamma-star}.
If we have moreover a section $s \in \Gamma(X, \mathcal{L})$, then
we obtain a map
\begin{equation}
\label{equation-module-invert-s}
\Gamma_*(X, \mathcal{L}, \mathcal{F})_{(s)}
\longrightarrow
\Gamma(X_s, \mathcal{F}|_{X_s})
\end{equation}
which sends $t/s^n$ where
$t \in \Gamma(X, \mathcal{F} \otimes_{\mathcal{O}_X} \mathcal{L}^{\otimes n})$
to $t|_{X_s} \otimes s|_{X_s}^{-n}$. This makes sense
because $X_s \subset X$ is by definition the open over which
$s$ has an inverse, see Modules, Lemma \ref{modules-lemma-s-open}.

\begin{lemma}
\label{lemma-invert-s-sections}
\begin{reference}
\cite[Section 6.8]{EGA1-second}
\end{reference}
Let $X$ be a scheme. Let $\mathcal{L}$ be an invertible sheaf on $X$.
Let $s \in \Gamma(X, \mathcal{L})$. Let $\mathcal{F}$ be a quasi-coherent
$\mathcal{O}_X$-module.
\begin{enumerate}
\item If $X$ is quasi-compact, then (\ref{equation-module-invert-s})
is injective, and
\item if $X$ is quasi-compact and quasi-separated, then
(\ref{equation-module-invert-s}) is an isomorphism.
\end{enumerate}
In particular, the canonical map
$$
\Gamma_*(X, \mathcal{L})_{(s)}
\longrightarrow
\Gamma(X_s, \mathcal{O}_X),\quad
a/s^n \longmapsto a \otimes s^{-n}
$$
is an isomorphism if $X$ is quasi-compact and quasi-separated.
\end{lemma}

\begin{proof}
Assume $X$ is quasi-compact. Choose a finite affine open covering
$X = U_1 \cup \ldots \cup U_m$ with $U_j$ affine and
$\mathcal{L}|_{U_j} \cong \mathcal{O}_{U_j}$. Via this isomorphism,
the image $s|_{U_j}$ corresponds to some
$f_j \in \Gamma(U_j, \mathcal{O}_{U_j})$. Then
$X_s \cap U_j = D(f_j)$.

\medskip\noindent
Proof of (1). Let $t/s^n$ be an element in the kernel of
(\ref{equation-module-invert-s}). Then $t|_{X_s} = 0$.
Hence $(t|_{U_j})|_{D(f_j)} = 0$. By
Lemma \ref{lemma-invert-f-sections} we conclude that
$f_j^{e_j} t|_{U_j} = 0$ for some
$e_j \geq 0$. Let $e = \max(e_j)$. Then we see that $t \otimes s^e$
restricts to zero on $U_j$ for all $j$, hence is zero. Since $t/s^n$
is equal to $t \otimes s^e/s^{n + e}$ in
$\Gamma_*(X, \mathcal{L}, \mathcal{F})_{(s)}$ we conclude that $t/s^n = 0$
as desired.

\medskip\noindent
Proof of (2). Assume $X$ is quasi-compact and quasi-separated.
Then $U_j \cap U_{j'}$ is quasi-compact for all pairs $j, j'$, see
Schemes, Lemma \ref{schemes-lemma-characterize-quasi-separated}.
By part (1) we know (\ref{equation-module-invert-s}) is injective.
Let $t' \in \Gamma(X_s, \mathcal{F}|_{X_s})$. For every $j$, there exist an
integer $e_j \geq 0$ and $t'_j \in \Gamma(U_j, \mathcal{F}|_{U_j})$ such that
$t'|_{D(f_j)}$ corresponds to $t'_j/f_j^{e_j}$
via the isomorphism of Lemma \ref{lemma-invert-f-sections}.
Set $e = \max(e_j)$ and
$$
t_j = f_j^{e - e_j} t'_j \otimes q_j^e \in
\Gamma(U_j,
(\mathcal{F} \otimes_{\mathcal{O}_X} \mathcal{L}^{\otimes e})|_{U_j})
$$
where $q_j \in \Gamma(U_j, \mathcal{L}|_{U_j})$ is the trivializing
section coming from the isomorphism
$\mathcal{L}|_{U_j} \cong \mathcal{O}_{U_j}$. In particular we have
$s|_{U_j} = f_j q_j$. Using this a calculation shows that
$t_j|_{U_j \cap U_{j'}}$ and $t_{j'}|_{U_j \cap U_{j'}}$
map to the same section of $\mathcal{F}$ over $U_j \cap U_{j'} \cap X_s$.
By quasi-compactness of $U_j \cap U_{j'}$ and part (1) there exists an
integer $e' \geq 0$ such that
$$
t_j|_{U_j \cap U_{j'}} \otimes s^{e'}|_{U_j \cap U_{j'}} =
t_{j'}|_{U_j \cap U_{j'}} \otimes s^{e'}|_{U_j \cap U_{j'}}
$$
as sections of $\mathcal{F} \otimes \mathcal{L}^{\otimes e + e'}$ over
$U_j \cap U_{j'}$. We may choose the same $e'$ to work for all pairs
$j, j'$. Then the sheaf conditions implies there is a section
$t \in \Gamma(X, \mathcal{F} \otimes \mathcal{L}^{\otimes e + e'})$
whose restriction to $U_j$ is $t_j \otimes s^{e'}|_{U_j}$.
A simple computation shows that $t/s^{e + e'}$ maps to $t'$
as desired.
\end{proof}

\noindent
Let $X$ be a scheme. Let $\mathcal{L}$ be an invertible $\mathcal{O}_X$-module.
Let $\mathcal{F}$ and $\mathcal{G}$ be quasi-coherent $\mathcal{O}_X$-modules.
Consider the graded $\Gamma_*(X, \mathcal{L})$-module
$$
M = \bigoplus\nolimits_{n \in \mathbf{Z}} \Hom_{\mathcal{O}_X}(\mathcal{F},
\mathcal{G} \otimes_{\mathcal{O}_X} \mathcal{L}^{\otimes n})
$$
Next, let $s \in \Gamma(X, \mathcal{L})$ be a section. Then there is a
canonical map
\begin{equation}
\label{equation-hom-invert-s}
M_{(s)} \longrightarrow
\Hom_{\mathcal{O}_{X_s}}(\mathcal{F}|_{X_s}, \mathcal{G}|_{X_s})
\end{equation}
which sends $\alpha/s^n$ to the map $\alpha|_{X_s} \otimes s|_{X_s}^{-n}$.
The following lemma, combined with
Lemma \ref{lemma-extend-finite-presentation},
says roughly that, if $X$ is quasi-compact and quasi-separated,
the category of finitely presented $\mathcal{O}_{X_s}$-modules
is the category of finitely presented $\mathcal{O}_X$-modules
with the multiplicative system of maps
$s^n: \mathcal{F} \to
\mathcal{F} \otimes_{\mathcal{O}_X} \mathcal{L}^{\otimes n}$ inverted.

\begin{lemma}
\label{lemma-section-maps-backwards}
Let $X$ be a scheme. Let $\mathcal{L}$ be an invertible $\mathcal{O}_X$-module.
Let $s \in \Gamma(X, \mathcal{L})$ be a section.
Let $\mathcal{F}$, $\mathcal{G}$ be quasi-coherent $\mathcal{O}_X$-modules.
\begin{enumerate}
\item If $X$ is quasi-compact and $\mathcal{F}$ is of finite type,
then (\ref{equation-hom-invert-s}) is injective, and
\item if $X$ is quasi-compact and quasi-separated and $\mathcal{F}$
is of finite presentation, then
(\ref{equation-hom-invert-s})
is bijective.
\end{enumerate}
\end{lemma}

\begin{proof}
We first prove the lemma in case $X = \Spec(A)$ is affine
and $\mathcal{L} = \mathcal{O}_X$. In this case $s$ corresponds
to an element $f \in A$. Say
$\mathcal{F} = \widetilde{M}$ and $\mathcal{G} = \widetilde{N}$
for some $A$-modules $M$ and $N$. Then the lemma translates
(via Lemmas \ref{lemma-finite-type-module} and
\ref{lemma-finite-presentation-module}) into
the following algebra statements
\begin{enumerate}
\item If $M$ is a finite $A$-module and $\varphi : M \to N$ is
an $A$-module map such that the induced map $M_f \to N_f$ is zero,
then $f^n\varphi = 0$ for some $n$.
\item If $M$ is a finitely presented $A$-module, then
$\Hom_A(M, N)_f = \Hom_{A_f}(M_f, N_f)$.
\end{enumerate}
The second statement is
Algebra, Lemma \ref{algebra-lemma-hom-from-finitely-presented} and we omit
the proof of the first statement.

\medskip\noindent
Next, we prove (1) for general $X$.
Assume $X$ is quasi-compact and hoose a finite affine open covering
$X = U_1 \cup \ldots \cup U_m$ with $U_j$ affine and
$\mathcal{L}|_{U_j} \cong \mathcal{O}_{U_j}$. Via this isomorphism,
the image $s|_{U_j}$ corresponds to some
$f_j \in \Gamma(U_j, \mathcal{O}_{U_j})$. Then
$X_s \cap U_j = D(f_j)$.
Let $\alpha/s^n$ be an element in the kernel of
(\ref{equation-hom-invert-s}). Then $\alpha|_{X_s} = 0$.
Hence $(\alpha|_{U_j})|_{D(f_j)} = 0$. By the affine case treated above
we conclude that $f_j^{e_j} \alpha|_{U_j} = 0$ for some
$e_j \geq 0$. Let $e = \max(e_j)$. Then we see that $\alpha \otimes s^e$
restricts to zero on $U_j$ for all $j$, hence is zero. Since $\alpha/s^n$
is equal to $\alpha \otimes s^e/s^{n + e}$ in $M_{(s)}$ we conclude that
$\alpha/s^n = 0$ as desired.

\medskip\noindent
Proof of (2). Since $\mathcal{F}$ is of finite presentation, the
sheaf $\SheafHom_{\mathcal{O}_X}(\mathcal{F}, \mathcal{G})$ is
quasi-coherent, see Schemes, Section \ref{schemes-section-quasi-coherent}.
Moreover, it is clear that
$$
\SheafHom_{\mathcal{O}_X}(\mathcal{F},
\mathcal{G} \otimes_{\mathcal{O}_X} \mathcal{L}^{\otimes n}) =
\SheafHom_{\mathcal{O}_X}(\mathcal{F}, \mathcal{G})
\otimes_{\mathcal{O}_X} \mathcal{L}^{\otimes n}
$$
for all $n$. Hence in this case the statement follows from
Lemma \ref{lemma-invert-s-sections} applied to
$\SheafHom_{\mathcal{O}_X}(\mathcal{F}, \mathcal{G})$.
\end{proof}





